\documentclass[11pt]{article}
\usepackage{../homework}

\profname{Dr. Stanislaw Szarek}
\classcode{MATH 423 Measure Theory}
\assignnum{Assignment 05}
\duedate{28 September 2026}

\begin{document}
All problems from Folland, \textit{Real Analysis} (2e).

\begin{problem}{1.5.30}
If \( E \in \mathcal{L} \) and \( m(E) > 0 \), for any \(\alpha < 1\) there is an open interval \( I \) such that \( m(E \cap I) > \alpha m(I) \).
\end{problem}
\begin{proof}
  We can reduce to the finite measure case; write \( E_n = E \cap [-n, n] \), and by continuity from below \( m(E_n) \rightarrow m(E) = \infty \).
  Then there is some \( N \in \mathbb{N} \) for which \( m(E_n) > 0 \) for all \( n \geq N \).
  Then the interval we find for \( E_n \) holds for all \( n \geq N \) and in particular for \( E \).
  Hence, assume \( m(E) < \infty \), and consider a open cover \( \{I_j\} \) of intervals such that \( E \subset \bigcup_j I_{j} \) and \( m(E) \approx \sum\limits_{j} m(I_j) \), that is, \( \sum\limits_{j} m(I_j) < m(E) + \epsilon \).
  Now assume, indirectly, that there exists some \( \alpha < 1 \) for which any open interval satisfies \( m(E \cap I) \leq \alpha m(I) \), and fix this \(\alpha\).
  Write the cover as \( E = E \cap \bigcup_{j} I_{j} \), and apply countable sub-additivity to get
  \begin{equation}
  \label{eq:2}
  m(E) \leq \sum\limits_{j} m(E \cap I_j) \leq \sum\limits_j \alpha m(I_j) = \alpha \sum\limits_j m(I_j) < \alpha(m(E) + \epsilon).
\end{equation}

We have that \( m(E) < \alpha(m(E) + \epsilon) \).
Equivalently, \( m(E) < \epsilon \frac{\alpha}{1-\alpha} \).
Since \( m(E) = \inf \sum\limits_j m(I_j) \) where the inf is taken over open covers of \( E \), \(\epsilon\) can be chosen arbitrarily.
In particular, choosing \( 0 < \epsilon < \frac{1-\alpha}{\alpha}m(E) \) immediately gives a contradiction.
\end{proof}

\begin{problem}{1.5.31}
  If \( E \in \mathcal{L} \) and \( m(E) > 0 \), the set \( E - E = \{x - y \mid x,y \in E\} \) contains an interval centered at \( 0 \).
  (If \( I \) is as in Exercise 30 with \(\alpha > \frac{3}{4}\), then \( E - E \) contains \( (-\frac{1}{2}m(I), \frac{1}{2}m(I)) \)).
\end{problem}
\begin{proof}
  We want to show that for some \(R\), \( (-R, R) \subset E - E \).
  The condition \(s \in E - E\) means that there exist \( x, y \in E \) such that \(s = x - y \) or \( x = y + s \).
  In other words, \( x \in (s + E) \) and in particular \( x \in E \cap (s + E) \).
  We now have an equivalent formulation of the claim: for all \( s \in (-R, R) \), \( s \in E - E \iff E \cap (s + E) \neq \emptyset \).

  From Problem 1.5.30, we see that for any \(\alpha < 1\), we have some \( I \) for which \( m(E \cap I) > \alpha m(I) \).
  Write \( I = (p, q) \), and fix \( |s| < \frac{1}{2}m(I) \).
  Consider now the sets \( E \cap I \) and \( s + (E \cap I) \).
  By the translation invariance of the Lebesgue measure,
  \begin{equation}
  \label{eq:5}
  m(E \cap I) = m(s + (E \cap I)) > \alpha m(I).
\end{equation}
Now, both \( (E \cap I), (s + (E \cap I)) \subset I \cup (s + I) \), which has \( m(I \cup (s + I)) = m(I) + |s| \).
Moreover, since \( (E \cap I) \subset E \), \( (E \cap I) \cap (s + (E \cap I)) \subset E \cap (s + E) \).

What remains is to show that \( (E \cap I) \cap (s + (E \cap I)) \neq \emptyset \).
Suppose the contrary, so that the two sets are disjoint.
From \cref{eq:5} and additivity of \( m \),
\begin{equation}
\label{eq:6}
m((E \cap I) \cup (s + (E \cap I))) = m(E \cap I) + m(s + (E \cap I)) = 2m(E \cap I) > 2\alpha m(I).
\end{equation}
Let \( \alpha > \frac{3}{4} \), \( |s| < \frac{1}{2}m(I) \).
We also have the bound that
\begin{equation}
\label{eq:8}
m((E \cap I) \cup (s + (E \cap I))) \leq m(I \cup (s + I)) = m(I) + |s| < \frac{3}{2}m(I).
\end{equation}

Comparing \cref{eq:6,eq:8} shows that
\begin{equation}
\label{eq:9}
2\alpha m(I) < m((E \cap I) \cup (s + (E \cap I))) < \frac{3}{2}m(I).
\end{equation}
The left hand bound is then \( 2\alpha m(I) > \frac{3}{2}m(I) \), which contradicts the right hand bound.
Therefore, \( (E \cap I) \cap (s + (E \cap I)) \neq \emptyset \), and our claim now follows.

\end{proof}

\begin{problem}{2.1.4}
If \( f: X \rightarrow \mathbb{R} \) and \( f^{-1}((r, \infty)) \in \mathcal{M} \) for each \( r \in \mathbb{Q} \), then \( f \) is measurable.
\end{problem}
\begin{proof}
  By definition, \( f \) is measurable if for any Borel set \( B \in \mathcal{B}_{\mathbb{R}} \), the inverse image \( f^{-1}(B) \in \mathcal{M} \) is measurable.
  Now, we know that (by Proposition 1.2) the Borel sigma algebra is generated by the open rays \( (a, \infty) \).
  These open rays can themselves be generated by rational rays; since \( \mathbb{Q} \) is dense in \( \mathbb{R} \),  for any \( a \in \mathbb{R} \) there is some decreasing sequence \( (q_n) \subset \mathbb{Q} \) for which \( q_n \searrow a \).
  Then,
  \begin{equation}
  \label{eq:3}
  (a,\infty) = \bigcup_{n=1}^{\infty} (q_n, \infty).
\end{equation}

All of this is to say, the Borel sigma algebra can be generated by these rays with a rational endpoint.
Explicitly, let \( \mathcal{E} \) be the set of rays with rational endpoint, and \( \mathcal{E}_5 \) be the set of open rays (which generates \( \mathcal{B}_{\mathbb{R}} \)); we have shown that \( \mathcal{E}_5 \subset \mathcal{M}(\mathcal{E})\) so that by Lemma 1.1, \( \mathcal{B}_{\mathbb{R}} = \mathcal{M}(\mathcal{E}) \).

By Proposition 2.1, if \( \mathcal{B}_{\mathbb{R}} \) is generated by \( \mathcal{E} \), then \( f \) is measurable if and only if \( f^{-1}(E) \in \mathcal{M} \) for every \( E \in \mathcal{E} \).
That is precisely the assumption in our claim, so we are done.
\end{proof}

\begin{problem}{2.1.6}
  The supremum of an uncountable family of measurable \( \mathbb{R} \)-valued functions on \( X \) can fail to be measurable (unless the \(\sigma\)-algebra \( \mathcal{M} \) is very special).
\end{problem}
\begin{proof}
  Let \( E \subset X \) be a non-measurable, uncountable set.
  For each \( x \in E \), define \( f_x = \chi_{\{x\}} \).
  Then, for \( y \in X \),
  \begin{equation}
  \label{eq:4}
  \sup_{x \in E} f_x(y) = \begin{cases}
    1, &y \in E \\
    0, &y \not\in E
  \end{cases} = \chi_E(y).
\end{equation}

Assuming singletons are measurable in \( X \) (take, for example, \( (X, \mathcal{M}) = (\mathbb{R}, \mathcal{L}) \)), each of \( \chi_{\{x\}} \) is measurable.
However, since \(\chi^{-1}_E({1}) = E\) which is not measurable in \( X \), we have a measurable family whose supremum is non-measurable.
\end{proof}

\begin{problem}{2.1.8}
If \( f: \mathbb{R} \rightarrow \mathbb{R} \) is monotone, then \( f \) is Borel measurable.
\end{problem}
\begin{proof}
  Suppose that \( f \) is non-decreasing, so that if \( x_1 < x_2 \), \( f(x_1) \leq f(x_{2}) \).
  We want to show that for any \( U \in \mathcal{B}_{\mathbb{R}} \), \( f^{-1}(U) \in \mathcal{B}_{\mathbb{R}} \).
  By Proposition 2.1 it suffices to show this holds for a generating family of \( \mathcal{B}_{\mathbb{R}} \), in particular we will use the open rays \( \mathcal{E}_5 \).
  Fix some \( a \in \mathbb{R} \) and consider \( f^{-1}((a, \infty)) = \{x \in \mathbb{R} \mid f(x) > a\} \).
  Since \( f \) is non-decreasing, if \( x \in f^{-1}((a, \infty)) \), then for any
  \begin{equation}
    \label{eq:monotonepreimg}
    y > x, \quad \text{ we have } y \in f^{-1}((a, \infty)).
  \end{equation}

  Defining \( A \coloneq f^{-1}((a,\infty)) \), we claim the only forms of this set can be \(\emptyset, \mathbb{R}, (b, \infty) \), or \([b, \infty)\), for some \( b \in \mathbb{R} \).
  The first two are obvious: if \( f \) is bounded from above by \( M < a \), then \( A = \emptyset \); on the other hand, if \( f \) is bounded from below by \( M > a \), then \( A = \mathbb{R} \).
  Now, let \( b \coloneq \inf A \).
  The condition imposed by the monontonicity, \cref{eq:monotonepreimg}, means that \( A \) is a ray.
  Then either \( A \) takes on the value of \( \inf A \), so \( A = [b, \infty) \), or it doesn't so that \( A = (b, \infty) \).

  All of these possible inverse images are Borel sets, so our claim follows.
  A symmetric argument with lower rays \( (-\infty, a) \) shows the inverse images for a non-increasing \( f \) are \( \emptyset, \mathbb{R}, (-\infty, b], (-\infty, b) \).
\end{proof}
\end{document}
